\section*{Resumen}
\addcontentsline{toc}{section}{Resumen}

El Sistema Inteligente de Gesti\'on de Aulas (SIGA) es un sistema IoT + IA propuesto para la Facultad de Ciencias de la Computaci\'on de la UTEQ, orientado a monitorear condiciones ambientales y ocupaci\'on de aulas, controlar remotamente proyectores y climatizaci\'on, generar alertas ante anomal\'ias, gestionar solicitudes de mantenimiento y automatizar el ahorro energ\'etico. El presente informe cierra el Proyecto Integrador de la asignatura Ingenier\'ia de Requerimientos (ISR-401) mediante la Unidad V: auditor\'ia de calidad del ERS con seis m\'etricas derivadas de ISO/IEC 25010:2023, verificaci\'on de trazabilidad end-to-end, especificaci\'on de requisitos para dos componentes de inteligencia artificial (predicci\'on de fallas y an\'alisis de consumo energ\'etico), y defensa t\'ecnica ante el tribunal docente. La auditor\'ia realizada sobre el documento real (primera inspecci\'on formal registrada del proyecto) encontr\'o que cuatro de las seis m\'etricas de calidad no alcanzan el valor de referencia del curso, con una contradicci\'on cr\'itica de priorizaci\'on identificada entre el documento narrativo y el archivo de priorizaci\'on vigente. Se declaran expl\'icitamente los hallazgos, las acciones de mejora aplicadas y las pendientes, sin maquillar cifras.

\textbf{Palabras clave:} ingenier\'ia de requisitos, IoT, gesti\'on de aulas, calidad de requisitos, trazabilidad, inteligencia artificial.

\bigskip
\textbf{Abstract}

The Intelligent Classroom Management System (SIGA) is an IoT + AI system proposed for the Faculty of Computer Science at UTEQ, aimed at monitoring environmental conditions and room occupancy, remotely controlling projectors and air conditioning, generating anomaly alerts, managing maintenance requests, and automating energy savings. This report closes the Integrative Project of the Requirements Engineering course (ISR-401) through Unit V: a quality audit of the SRS using six metrics derived from ISO/IEC 25010:2023, end-to-end traceability verification, requirements specification for two artificial-intelligence components (failure prediction and energy-consumption analysis), and a technical defense before the faculty panel. The audit performed on the actual document (the first formal inspection ever recorded for this project) found that four of the six quality metrics fail to reach the course's reference threshold, with a critical prioritization contradiction identified between the narrative document and the current prioritization file. Findings, applied improvements, and open items are stated explicitly, without inflating figures.

\textbf{Keywords:} requirements engineering, IoT, classroom management, requirements quality, traceability, artificial intelligence.
